\paragraph{Maximisation bias and double estimators.} The positive bias of the maximum of noisy estimates is classical in decision analysis~\cite{smith2006optimizer}. van Hasselt (2010) introduced Double Q-learning and the trap MDP of Sutton and Barto's textbook that we use (neither has an arXiv or Crossref record, so they are not in our reference list). Double DQN~\cite{hasselt2015deep} transfers the idea to deep networks using the target network as the second estimator. Weighted Double Q-learning~\cite{zhang2017weighted}, whose abstract notes that double Q-learning sometimes underestimates action values, interpolates between the single and the double estimator with a data-dependent weight; we reimplement it as a baseline. Maxmin Q-learning~\cite{lan2020maxmin} bootstraps from the minimum over several estimators, with the number of estimators controlling the bias, and its abstract highlights that the effect of overestimation on learning efficiency is environment dependent, a point our random-MDP results echo.

\paragraph{Underestimation and bias control.} Ren et al.~\cite{ren2021estimation} argue that the underestimation of double Q-learning can lead to non-optimal fixed points under an approximate Bellman operator. Cini et al.~\cite{cini2020deep} study a weighted estimator in deep RL, Sabry and Khalifa~\cite{sabry2019reduction} use dropout to reduce variance and overestimation, and Kuznetsov et al.~\cite{kuznetsov2021automating} propose automatic selection of bias-control hyper-parameters for algorithms that include Maxmin Q-learning. Wagenbach and Sabatelli~\cite{wagenbach2022factors} study how learning rate, discount and reward signal influence Q-learning's overestimation in settings that require function approximation, which motivates our step-size and noise sweeps in a setting where exact values are available. TD3~\cite{fujimoto2018addressing} and Rainbow~\cite{hessel2017rainbow} are examples of systems in which such corrections are standard components.

\paragraph{Gap.} Most of the above evaluates new estimators, often with function approximation. We instead ask a narrow, controlled question with exact ground truth: in tabular settings with matched hyper-parameters, how do the standard estimators compare, and which confounds explain the differences? We do not propose a new estimator.
